\section{Preliminaries and Problem Formulation}
\label{sec:problem-formulation}

Let $\pi_{\theta_0}$ be a fixed base policy, $\mathcal{B}$ a source dataset with
a verifier $r_{\mathcal{B}}(x,y)$, and $\mathcal{A}$ a target benchmark excluded
from RLVR training.  Training method $q\in\{\mathrm{SFT},\mathrm{RSFT},
\mathrm{RLVR}\}$ on $\mathcal{B}$ produces policy
$\pi_{\theta_{q,\mathcal{B}}}$.  We define directed transfer as
\begin{equation}
    T^{q}_{\mathcal{B}\rightarrow\mathcal{A}}
    = \operatorname{Score}(\pi_{\theta_{q,\mathcal{B}}},\mathcal{A})
      - \operatorname{Score}(\pi_{\theta_0},\mathcal{A}).
    \label{eq:directed-transfer}
\end{equation}
A positive value denotes positive transfer, a negative value denotes negative
transfer, and uncertainty intervals that include zero are treated as
inconclusive rather than evidence of transfer.

\paragraph{Information-access protocol.}
The deployable setting permits source prompts, source answers and rewards, and
a disjoint unlabeled target development sample $\mathcal{A}^{u}$.  It prohibits
target answers, target rewards, and target-test feedback during method design,
data selection, hyperparameter tuning, and checkpoint selection.  Once target
test results have influenced a design decision, that target is considered a
development benchmark and cannot serve as a confirmatory test for the revised
method.

\paragraph{Research questions.}
We ask (i) when RLVR transfers across benchmarks, (ii) whether directed
reasoning coverage predicts this transfer beyond semantic similarity and task
difficulty, (iii) whether observed gains reflect empirical support expansion
or redistribution, and (iv) whether the predictive signal can improve
fixed-budget RLVR data selection.

